% chapters/07/section-05.tex
\section{Program Tabrakan Lengkap}
\label{sec:program-tabrakan}

Bagian ini menyatukan seluruh unsur menjadi satu prototipe permainan yang dapat
dimainkan: pemain bergerak, jatuh, melompat, mendarat di bongkahan, terhalang
oleh sisinya, mengumpulkan kepingan, dan kehilangan nyawa ketika menyentuh
rintangan.

\subsection{Kode Lengkap}

\kodepython{python/ch07_tabrakan.py}{Permainan dengan tabrakan per sumbu, kepingan, dan rintangan}{lst:ch07-tabrakan}

\subsection{Tiga Pemakaian Tabrakan di Dalamnya}

Ketiga perintah pemeriksaan tabrakan dipakai untuk keperluan
yang berbeda, sesuai pertimbangan pada Tabel~\ref{tab:perintah-tabrakan}.

\begin{description}
  \item[Bongkahan, dengan \code{colliderect} dan penyelesaian per sumbu.]
        Diperiksa di dalam \code{Player.update}. Pemeriksaan ini bukan sekadar
        mengetahui, melainkan memperbaiki posisi dan kecepatan. Hasilnya
        adalah kemampuan mendarat di atas bongkahan dan terhalang oleh
        sisinya, yaitu indikator \ind{b}.
  \item[Kepingan, dengan \code{spritecollide} bernilai benar.]
        Diperiksa pada tahap pembaruan. Kepingan yang tersentuh langsung
        dikeluarkan dari wadahnya, dan jumlahnya dihitung dari panjang daftar
        yang dikembalikan.
  \item[Rintangan, dengan \code{spritecollide} bernilai salah.]
        Diperiksa dengan cara yang sama, tetapi rintangannya tidak
        dikeluarkan, sebab rintangan harus tetap ada di tempatnya. Sentuhan
        mengurangi satu nyawa dan mengembalikan pemain ke titik awal.
\end{description}

\subsection{Pemeriksaan yang Dijalankan Sepanjang Waktu}

Program ini memeriksa kebenaran penanganan tabrakannya sendiri pada setiap
bingkai, sebagaimana dianjurkan pada Bagian~\ref{sec:per-sumbu}. Penghitung
tumpang tindih ditampilkan pada terminal ketika program berakhir.

\begin{kodekeluaran}
Bingkai dijalankan       : 239
Kepingan terkumpul       : 1 dari 5
Skor                     : 10
Tersentuh rintangan      : 3
Bingkai bertumpang tindih: 0
\end{kodekeluaran}

Baris terakhir adalah yang paling penting: nilai nol menunjukkan bahwa pada
seluruh bingkai yang dijalankan, pemain tidak pernah berada di dalam
bongkahan. Dengan kata lain, penanganan tabrakannya benar. Angka itu tetap nol
walaupun pemain bergerak cepat, melompat, dan terhalang berkali-kali.

\begin{catatabox}
Program itu dijalankan dengan pilihan \code{--auto-run} yang
menggerakkan pemain secara otomatis. Kemampuan menjalankan program tanpa
sentuhan manusia semacam itu berguna untuk pengujian: Anda dapat menjalankan
ratusan bingkai dalam beberapa detik, dan memeriksa kebenaran penanganan
tabrakan tanpa harus memainkannya sendiri.
\end{catatabox}

\subsection{Percobaan yang Disarankan}

\begin{table}[htbp]
\centering
\caption{Percobaan pada program tabrakan}
\label{tab:percobaan-ch07}
{\setlength{\tabcolsep}{4pt}\footnotesize%
\begin{tabularx}{\textwidth}{@{}L{5.6cm} Y@{}}
\toprule
\headrow \textbf{Yang diubah dan caranya} & \textbf{Yang sebaiknya Anda amati} \\
\midrule
Menyatukan kedua sumbu, dengan memindahkan posisi pada kedua sumbu sekaligus
lalu memeriksa tabrakan sekali saja. &
Pemain yang menabrak sisi bongkahan terangkat ke atasnya, dan pemain yang
mendarat kadang tersangkut di udara. Kembalikan sesudah percobaan. \\
\addlinespace
Menghapus penolan kecepatan tegak sesudah mendarat. &
Pemain bergetar halus di permukaan bongkahan, sebab posisinya diperbaiki
berulang kali pada setiap bingkai. \\
\addlinespace
Memperbesar kecepatan pemain sampai di atas 1500 piksel per detik. &
Pemain dapat menembus bongkahan yang tipis, karena pada satu bingkai
perpindahannya melampaui tebal bongkahan. \\
\addlinespace
Menghapus pemeriksaan \code{count_overlaps}. &
Program tetap berjalan, tetapi Anda kehilangan alat untuk mengetahui apakah
penanganan tabrakannya masih benar. \\
\addlinespace
Mengubah nilai \code{dokill} menjadi salah pada pemeriksaan kepingan. &
Kepingan tidak hilang ketika tersentuh, dan skor bertambah setiap bingkai
selama pemain menyentuhnya. \\
\bottomrule
\end{tabularx}}
\end{table}

\begin{peringatanbox}
Percobaan pertama pada tabel tersebut adalah yang paling penting untuk dicoba.
Mengerjakan penanganan tabrakan dengan cara yang salah, lalu mengamati
gejalanya sendiri, jauh lebih berkesan daripada hanya membacanya. Catatlah
gejala yang Anda amati pada lembar kerja praktikum; gejala itulah yang menjadi
bahan analisis pada Bagian~\ref{sec:analisis-tabrakan}.
\end{peringatanbox}
